\scanpage{45}
\uline{Formes extérieures types}

\margin{$E_0=\operatorname{Ext}(G_0)$\\$H_0=\operatorname{ad}(G_0)$\\$\Gamma_0=E.H$}

$G_0$ un groupe \unsure{s. simple} type (sur $\mathbb Z$), $\Gamma_0=\underline{\operatorname{Aut}}_{\mathbb Z}(G_0)$, \deleted{$\unsure{E=\Gamma/\Gamma^0}$} ou \deleted{$\unsure{E_0}$}.
\inserted{$E_0$ \unsure{sous-groupe} normalisant $T_0\subset B_0$}
$S$ schéma de base

$G$ forme de $G_0$ sur $S$, correspond à un principal homogène \deleted{$\unsure{Q_G}$}\inserted{$\unsure{Q_0}$} $\Gamma_S$, \deleted{\unsure{sous}}
\deleted{$\unsure{Q_G/E_0}$}\inserted{$\rho=\rho_G$}, principal homogène associé de groupe $E$, \deleted{qui} est \unsure{noté} par abus de langage l'élément correspondant de $H^1(S,E_0)$.

\deleted{\unsure{le principal} $\unsure{G_0^\rho}$}
$G_1=G_0^\rho$ déduit de $G_0$ en tordant par $\rho$ \inserted{[\unsure{par l'interméd.} \unsure{via} : $\unsure{\rho:E_0}$]} : une telle \deleted{groupe} forme sera dite \uline{forme extérieure type} associée à $\rho$.
[On dit que $G$ est une forme extérieure type s'il est \uline{isom.} à $G_{\unsure{1|0}}^\rho$, \deleted{$\unsure{G}$} \unsure{on a nécessairement} $\rho\simeq\rho_G$]

Si on part de $\rho$, [sans $G$], \uline{la donnée d'un $G$ et d'un isom} $\rho\xrightarrow[\sim]{\varphi}\rho_G$ \uline{``équivaut''} \uline{à la donnée d'un fibré principal homogène sous} $H_1=\unsure{H_0^\rho}$, la classification des objets $G$ à isom près :
[un isom de $(G,\varphi)$ et $(G',\varphi')$ étant donné par un isom $\psi:G\longrightarrow G'$ tel que l'isom correspondant $\rho_G\xrightarrow{\sim}\rho_{G'}$ soit compatible avec les deux $\rho_G$ et $\rho_{G'}$ \unsure{avec} $\varphi$ $\varphi'$ ; \uline{un tel automorphisme est dit intérieur}] \deleted{\unsure{revient}} \deleted{à classifier} est donnée par $H^1(S,\unsure{H_1})$. On a

\margin{Distinguer entre les deux \unsure{versions} avec $G_0$ et $H_0$, et avec $G_1$ et $H_1$ \unsure{-----}}
\margin{\unsure{Relier ici} la notion d'isom. \uline{intérieur} d'un \unsure{groupe}.}

\diagram{\begin{tikzcd}[ampersand replacement=\&,column sep=large]
H^1(S,\unsure{H_1}) \arrow[r] \& H^1(S,\Gamma) = \begin{array}{l}\text{ens des classes d'isom}\\\text{des formes de}\\G\text{ sur }S\end{array}\\
\begin{array}{l}\text{``ens des classes à}\\\text{isom près de ``formes}\\\deleted{\text{\unsure{extérieures}}}\ \text{intérieures relatives}\\\text{à }\rho\text{''}\end{array} \&
\end{tikzcd}}
\margin{44}
